\naslov{Gibanje v polju centralne sile}

\begin{definicija}
  Sila $\vec{F} = \vec{F}(\mb{P})$ je \pojem{centralna}, če obstaja točka
  $\mb{O}$ (pol sile), da $\vec{F}$ deluje v smeri zveznice med $\mb{P}$ in
  $\mb{O}$, in da je njena velikost odvisna le od razdalje.
\end{definicija}

\begin{trditev}
  Konzervativna sila $\vec{F}$ je centralna natanko tedaj, ko obstaja pol, okoli
  katerega je vrtilna količina konstantna.
\end{trditev}

\begin{proof}
  Recimo, da je $\vec{F}$ centralna.
  Tedaj za vrtilno količino okoli $\mb{O}$, velja
  \begin{gather*}
	\vec{l}(\mb{O}, \mb{P}) = (\mb{P}-\mb{O}) \times m \dot{\mb{P}} \\
	\partial_t \vec{l} = \dot{\mb{P}} \times m \dot{\mb{P}}
	+ (\mb{P} - \mb{O}) \times m \ddot{\mb{P}}
	= (\mb{P} - \mb{O}) \times \vec{F}
  \end{gather*}
  Če je $\vec{F}$ centralna, je vzporedna $\mb{P} - \mb{O}$, in se izniči tudi
  drugi člen.

  Recimo, da je vrtilna količina konstantna.
  Po zgornjem izračunu $\dot{\vec{l}} = (\mb{P} - \mb{O}) \times \vec{F} = 0$,
  torej je $\vec{F}$ vzporedna zveznici.
  Pokazati moramo še, da je velikost sile odvisna le od razdalje.
  Po predpostavki je sila konzervativna, torej $\vec{F} = -\grad U$ in
  \[
	\vec{F} = - \left(
	  \partial_r U \vec{e}_r + \inv{r} \partial_\theta U \vec{e}_\theta + \inv{r
		\sin \theta} \partial_\varphi U \vec{e}_\varphi
	\right)
  \]
  v sferičnih koordinatah.
  Ker je sila vzporedna zveznici $\vec{e}_r$, mora veljati $U = U(r)$.
\end{proof}

\vprasanje{Definiraj centralno silo in jo karakteriziraj ob predpostavki, da je
  konzervativna. Dokaži karakterizacijo.}

\begin{trditev}
  Zvezna centralna sila je konzervativna.
\end{trditev}

\begin{proof}
  Definiramo
  \[
	U = \int_{\abs{\mb{P}_0 - \mb{O}}}^{\abs{\mb{P}- \mb{O}}} F(r) dr +
	U(\mb{P}_0).
  \]
\end{proof}

\vprasanje{Dokaži: zvezna centralna sila je konzervativna.}